\chapter*{\center \Large  Abstract}

We train/assess MADRL players on a 2vs2 competitive football setting in a MuJoCo physics environment utilizing PettingZoo interface to allow fully autonomous players to learn how to chase a ball, play defence, and work together to score a goal from raw input data and no manual rules. Our primary methodology uses PPO reinforcement learning with a shared agent parameterization, enabling N/A solution for continuous action/state spaces, and N/A solution for extremely sparse reward issues were addressed by implementing/using custom reward shaping wrapper that calculates instantaneous rewards based on physical metrics such as distance penalties for chasing the ball, and speed of players toward goal.

In addition, we evaluate tactical play via spatial analysis software which produces 2D pitch coverage heatmaps, team spatial distribution (spread), and team possession statistics. Training occurs via parallelized vectorized training runs on a single Google Colab T4 GPU. The resulting PPO agents exhibit excellent ball interception skills and team coordination/individual coordination. Additionally, the hexbin maps of the pitches indicate that teams divide the pitch naturally into offensive/defensive zones (by colour) in an approximate 50:50 possession distribution with 48\% (red) and 51\% (blue). Continuous reward wrappers and agent parameter sharing produce agents that have learned to overcome physics-based non-stationarity.


~\\[1cm]
\noindent
\textbf{Keywords:} Multi-Agent Reinforcement Learning, Proximal Policy Optimisation, Parameter Sharing, Reward Shaping, MuJoCo Physics, Soccer Analytics, Spatial Heatmaps, PettingZoo.
